\chapter{Asymptotic normality}
\label{chap:normality}

This chapter proves the second main theorem: refined $o_p(\sqrt n)$ rates for the
deviation between proportions and plug-in targets, and joint central limit
theorems for the estimators and the allocation proportions. As before, everything
is derived from the generic conditions, so it applies to every $\aRTS$ design via
Proposition~\ref{prop:aRTS_generic}. It corresponds to Theorem~3.2 and
Appendix~A.3 of the paper.

We keep the standing assumptions of Chapter~\ref{chap:consistency}: a RAR
procedure satisfying the generic conditions with some $\ell_{n,k}$, and
Conditions~\textbf{A}--\textbf{B}. We introduce the notation
\[
  G := \nabla\Target\big|_\Params, \qquad
  V := \mathrm{diag}\!\Big(\tfrac{V_1}{v_1}, \dots, \tfrac{V_K}{v_K}\Big), \qquad
  \Omega := \begin{pmatrix} GVG^\top & GVG^\top\\ GVG^\top & GVG^\top\end{pmatrix}.
\]

\section{A refined decomposition of \texorpdfstring{$U$}{U}}
\label{sec:norm_U}

\begin{lemma}[Additive decomposition of the $U$-increment]\label{lem:diff_U_decomp}
\lean{AlphaRAR.diff_U_decomp}\leanok
\uses{def:U}
For any non-decreasing sequence $(\ell_n)$ with $\ell_n\le n$, under
Conditions~\textbf{A}--\textbf{B}, for each $k$,
\[
  U_{n,k} - U_{\ell_n,k}
  = (n-\ell_n)\big[-(1-\alpha)v_k + o(1)\big]
    + (M_{n,k} - M_{\ell_n,k})
    + \ell_n\big(\hat{\rho}_{\ell_n,k} - \hat{\rho}_{n,k}\big).
\]
\end{lemma}

\begin{proof}\leanok
\uses{thm:LLN}
Expanding $U_{n,k}-U_{\ell_n,k}$ from Definition~\ref{def:U} and grouping the
target terms around $v_k$,
$U_{n,k}-U_{\ell_n,k}
 = (n-\ell_n)\big[-(1-\alpha)v_k - (\hat{\rho}_{n,k}-v_k)
   + \frac{\alpha}{n-\ell_n}\sum_{m=\ell_n}^{n-1}(\hat{\rho}_{m,k}-v_k)\big]
   + M_{n,k}-M_{\ell_n,k} + \ell_n(\hat{\rho}_{\ell_n,k}-\hat{\rho}_{n,k})$.
By Theorem~\ref{thm:LLN}, $\hat{\rho}_n\to v$, so both $\hat{\rho}_{n,k}-v_k$ and
the Cesàro average $\frac1{n-\ell_n}\sum_{m=\ell_n}^{n-1}(\hat{\rho}_{m,k}-v_k)$
are $o(1)$ (treating separately the bounded and diverging cases of $\ell_n$).

\emph{Formalization.} \texttt{AlphaRAR.diff\_U\_decomp}. The exact $M$-explicit identity is
\texttt{AlphaRAR.auxU\_sub}; the $o(1)$ remainder $\alpha\sum_{[\ell_n,n)}(\hat\rho_m-v)
-(n-\ell_n)(\hat\rho_n-v)$ is made rigorous as ``for every $\varepsilon>0$, eventually
$|\cdot|\le\varepsilon(n-\ell_n)$'', its first term controlled by the windowed Cesàro bound
\texttt{AlphaRAR.abs\_sum\_Ico\_sub\_le\_of\_tendsto} (valid for an arbitrary window, bounded or
diverging) and its second by $\hat\rho_n\to v$.
\end{proof}

\begin{lemma}[Control of $\ell_n(\hat{\rho}_{\ell_n}-\hat{\rho}_n)$]\label{lem:ell_rho_control}
\lean{AlphaRAR.ell_rho_control}\leanok
\uses{def:plugin_target, def:Q}
Under Conditions~\textbf{A}--\textbf{B} there is a constant $C\ge0$ with, for any
non-decreasing $(\ell_n)$, $\ell_n\le n$,
\[
  \big|\ell_n(\hat{\rho}_{\ell_n} - \hat{\rho}_n)\big|
  \le o(1)(n-\ell_n) + C\|Q_n - Q_{\ell_n}\| + o_p(\sqrt n).
\]
\end{lemma}

\begin{proof}\leanok
\uses{lem:theta_error_Q, lem:rho_rate, cor:mart_Op, lem:taylor_rho}
By Corollary~\ref{cor:mart_Op}, $Q_{n,k}=O_p(\sqrt n)$, hence
$\hat{\theta}_{n,k}-\theta_k = Q_{n,k}/N_{n,k} + o_p(N_{n,k}^{-1/2}) = O_p(n^{-1/2})$
using $N_{n,k}\sim v_k n$ (Lemma~\ref{lem:theta_error_Q}). The Taylor expansion of
$\Target$ gives $\hat{\rho}_n - v = G(\estParams_n-\Params)+O_p(1/n)$ and similarly
at $\ell_n$; subtracting,
$\ell_n(\hat{\rho}_{\ell_n}-\hat{\rho}_n)
 = \ell_n G(\estParams_{\ell_n}-\estParams_n) + O_p(1)$, so
$|\ell_n(\hat{\rho}_{\ell_n}-\hat{\rho}_n)| \le C\ell_n\|\estParams_{\ell_n}-\estParams_n\| + O_p(1)$.
Writing the estimator difference through $Q$ (splitting into the increment
$\sum_{j=\ell_n+1}^n$ and the reweighting of $\sum_{j\le\ell_n}$) and bounding
each piece using $N_{n,k}-N_{\ell_n,k}\le n-\ell_n$ and the SLLN, the three
resulting terms are respectively $C\|Q_n-Q_{\ell_n}\|$, $o(1)(n-\ell_n)$ and
$o_p(\sqrt n)$.

\emph{Formalization.} \texttt{AlphaRAR.ell\_rho\_control} produces the bound in the shape
$\ell_n|\hat\rho_{\ell_n,k}-\hat\rho_{n,k}| \le V_\rho(n-\ell_n)+W_\rho$ with $V_\rho=o_p(1)$,
$W_\rho=o_p(\sqrt n)$, which is what the drift-sign step consumes. The Condition-\textbf{B} input
enters as the single hypothesis \texttt{hlip} (the target is Lipschitz near $\Params$, so
$|\hat\rho_{\ell,k}-\hat\rho_{n,k}|\le L\sum_{k'}|\hat\theta_{\ell,k'}-\hat\theta_{n,k'}|$
eventually) --- mirroring how \texttt{clt\_rho} takes the derivative as a hypothesis. The
estimator-increment control $\ell|\hat\theta_{\ell,k'}-\hat\theta_{n,k'}|\le g_{k'}(n-\ell)+h_{k'}
|Q_{n,k'}-Q_{\ell,k'}|$ is \texttt{AlphaRAR.abs\_estimator\_diff\_le}; the coefficients are
$g_{k'}=o_p(1)$ (from the Doob-maximal bound \texttt{isBigOpOne\_sup'\_abs\_div\_sqrt} and the a.s.\
bound \texttt{isBigOpOne\_of\_ae\_bounded}) and the reweighting coefficient
$h_{k'}=\ell_n/(N_{n,k'}+1)$, which is only \emph{bounded in probability} ($O_p(1)$, converging a.s.\
to $1/v_{k'}$, again via \texttt{isBigOpOne\_of\_ae\_bounded}), not by a uniform constant; it
therefore enters the $o_p$ bookkeeping through $O_p\cdot o_p=o_p$
(\texttt{IsBigOpOne.mul\_littleOp}). The $Q$-increment control
$|Q_{n,k'}-Q_{\ell,k'}|\le Q_v(n-\ell)+Q_w$ is \texttt{qm\_increments\_resp}.
\end{proof}

\section{Sharp increment bounds}
\label{sec:norm_bounds}

\begin{lemma}[$Q$- and $M$-increments are $o_p(\sqrt n)$ plus drift]\label{lem:QM_increments}
\lean{AlphaRAR.qm_increments_resp, AlphaRAR.qm_increments_of_bdd,
AlphaRAR.norm_increment_le_vmaxSeq_wmaxSeq}\leanok
\uses{def:Q, def:M, def:hitting}
Under Conditions~\textbf{A}--\textbf{B}, for the hitting time $\ell_{n,k}$,
\[
  \|Q_n - Q_{\ell_{n,k}}\| \le o_p(1)(n-\ell_{n,k}) + o_p(\sqrt n),
  \qquad
  |M_{n,k} - M_{\ell_{n,k},k}| \le o_p(1)(n-\ell_{n,k}) + o_p(\sqrt n).
\]
\end{lemma}

\begin{proof}\leanok
\uses{lem:mart_maximal, lem:mart_maximal_dyadic, lem:mart_maximal_pi, lem:increment_control,
lem:expectation_to_O}
Pick $L_n\to\infty$ with $L_n = o(n)$. Lemma~\ref{lem:mart_maximal} and
Lemma~\ref{lem:mart_maximal_dyadic} bound the two maximal quantities
$\E[\max_{m\le L_n}\|Q_n-Q_{n-m}\|]\le 2\sqrt{L_n C_0}$ and
$\E[\max_{L_n\le m<n}\frac{\|Q_n-Q_{n-m}\|}{m}]\le 3\sqrt{C_0/L_n}$; by
Lemma~\ref{lem:expectation_to_O} these are $O_p(\sqrt{L_n})=o_p(\sqrt n)$ and
$O_p(1/\sqrt{L_n})=o_p(1)$ respectively. Plugging into the deterministic
inequality of Lemma~\ref{lem:increment_control} gives the bound for $Q$; the same
argument applied to $M$ gives the second bound.

\emph{Formalization.} The two maximal quantities are formalized as \texttt{wmaxSeq} and
\texttt{vmaxSeq} (with the concrete window $L_n=\lfloor\sqrt n\rfloor$), and the two $o_p$ statements
are \texttt{AlphaRAR.isLittleOpOne\_wmaxSeq\_div\_sqrt} ($o_p(\sqrt n)$) and
\texttt{AlphaRAR.isLittleOpOne\_vmaxSeq} ($o_p(1)$), proved from the \emph{vector} maximal bounds
(Lemma~\ref{lem:mart_maximal_pi}) and Markov's inequality. The deterministic increment control is
\texttt{AlphaRAR.norm\_increment\_le\_vmaxSeq\_wmaxSeq}. For $Q$ (the response martingale, a
square-integrable family) the $o_p$ statements are \texttt{AlphaRAR.qm\_increments\_resp}; for the
scalar assignment martingale $M$ (bounded increments $|\Delta M|\le1$) they are
\texttt{AlphaRAR.qm\_increments\_of\_bdd} (via the singleton family $\{M\}$).
\end{proof}

\begin{lemma}[Upper bound on the $U$-increment]\label{lem:U_increment_bound}
\lean{AlphaRAR.isLittleOpOne_max_of_decomp}
\uses{def:U, def:hitting}
Under Conditions~\textbf{A}--\textbf{B}, for every $k$,
\[
  U_{n,k} - U_{\ell_{n,k},k} \le o_p(\sqrt n),
  \qquad
  U_{n,k} - U_{\ell_{n,k},k} \le O(\sqrt{n\log\log n})\ \text{a.s.}
\]
The $o_p(\sqrt n)$ half is formalized (generically) as \texttt{AlphaRAR.isLittleOpOne\_max\_of\_decomp};
the a.s.\ $O(\sqrt{n\log\log n})$ half is not yet formalized.
\end{lemma}

\begin{proof}
\uses{lem:diff_U_decomp, lem:ell_rho_control, lem:QM_increments, thm:LLN, thm:lil_bounded}
From Lemma~\ref{lem:ell_rho_control} and the $Q$-bound of
Lemma~\ref{lem:QM_increments},
$|\ell_{n,k}(\hat{\rho}_{\ell_{n,k}}-\hat{\rho}_{n,k})| \le o_p(1)(n-\ell_{n,k}) + o_p(\sqrt n)$.
Substituting this and the $M$-bound into the decomposition of
Lemma~\ref{lem:diff_U_decomp},
$U_{n,k}-U_{\ell_{n,k},k} = (n-\ell_{n,k})[-(1-\alpha)v_k+o_p(1)] + o_p(\sqrt n)
 \le o_p(\sqrt n)$, since the drift coefficient is $\le0$ and $n-\ell_{n,k}\ge0$.
For the a.s.\ bound, Theorem~\ref{thm:LLN} gives $n(\hat{\rho}_n-v)=O(\sqrt{n\log\log n})$,
hence $\ell_{n,k}(\hat{\rho}_{\ell_{n,k}}-\hat{\rho}_{n,k})=O(\sqrt{n\log\log n})$,
and the LIL for the martingale $M$ gives $M_{n,k}=O(\sqrt{n\log\log n})$; the same
substitution yields the second inequality.
\end{proof}

\begin{lemma}[Deviation between proportions and plug-in target]\label{lem:prop_dev}
\lean{AlphaRAR.aRTS_prop_dev, AlphaRAR.prop_dev}
\uses{def:generic_cond, def:counts, def:plugin_target}
Under the generic conditions and Conditions~\textbf{A}--\textbf{B}, for every $k$,
\begin{align*}
  |N_{n,k} - n\hat{\rho}_{n,k}| &= o_p(\sqrt n), &
  |N_{n,k} - n\hat{\rho}_{n,k}| &= O(\sqrt{n\log\log n})\ \text{a.s.}, \\
  N_{n,k} - n v_k &= O(\sqrt{n\log\log n})\ \text{a.s.}
\end{align*}
\emph{Formalization.} The generic $o_p(\sqrt n)$ assembly is \texttt{AlphaRAR.prop\_dev}: for a finite
family with $\sum_k \mathrm{Dev}_k = 0$ (\texttt{sum\_count\_sub\_smul\_eq\_zero}), the generic
inequality \eqref{eq:generic_ineq}, and the $U$-increment decomposition (drift $-(1-\alpha)v_k$, the
$M$- and $\rho$-terms increment-controlled by \texttt{qm\_increments\_of\_bdd} and
\texttt{ell\_rho\_control}, and an $o_p(1)$ perturbation from \texttt{diff\_U\_decomp}), it yields
$\mathrm{Dev}_k=o_p(\sqrt n)$ via the drift-sign lemma
\texttt{isLittleOpOne\_max\_div\_sqrt\_of\_drift} and the simplex reverse step
\texttt{isLittleOpOne\_dev\_of\_sum\_zero}. The concrete paper-form statement
$|N_{n,k}-n\hat\rho_{n,k}|=o_p(\sqrt n)$ is \texttt{AlphaRAR.aRTS\_prop\_dev}, the full, self-contained
$\aRTS$ instantiation: from the $\aRTS$ design bundle (an \texttt{IsAlgEnvSeq} sequence, Condition
\textbf{A} $Y\in L^2$, a continuous simplex-valued target $\Target$, the algorithm-level predicate
\texttt{IsARTS} --- whence the throttle via \texttt{throttle\_of\_isARTS} ---, $\alpha\in[0,1)$,
Condition \textbf{B} ($\Target$ Lipschitz near $\Params$ and non-sparsity $0<v_k$), and the
$\mathrm{thm{:}LLN}$ consistencies $N_{n,k}/n\to v_k$ and $\estParams_n\to\Params$), it discharges
\emph{every} hypothesis of \texttt{prop\_dev} internally: the generic inequality via
\texttt{generic\_ineq\_of\_hitting} (a.e., from the throttle), the smallness via
\texttt{preliminary\_small}, the $M$-increment via the seam $M=\texttt{assignMart}$
(\texttt{assignMart\_eq\_assignMG}) and \texttt{qm\_increments\_of\_bdd}, the \texttt{diff\_U\_decomp}
perturbation shown $o_p$ (a.e.\ $\to0$ by $\hat\rho_n\to v$), and the plug-in-target increment control
(Lemma~\ref{lem:ell_rho_control}) via \texttt{ell\_rho\_control}, whose $g$- and $h$-coefficients are
the random-hitting-time bounds \texttt{aRTS\_g\_littleOp} ($O_p\cdot o_p=o_p$, Doob running-max plus
$\sqrt n/(N_n+1)\to0$) and \texttt{aRTS\_h\_bigOp} (a.s.\ bounded). All random-hitting-time terms are
made measurable by \texttt{measurable\_eval\_of\_le}. Only the two a.s.\ $O(\sqrt{n\log\log n})$
statements remain unformalized.
\end{lemma}

\begin{proof}
\uses{lem:U_increment_bound, lem:counts_sum, thm:LLN}
The generic inequality \eqref{eq:generic_ineq} together with
\eqref{eq:generic_small} and Lemma~\ref{lem:U_increment_bound} gives the one-sided
bounds $N_{n,k}-n\hat{\rho}_{n,k} \le o_p(\sqrt n)$ and $\le O(\sqrt{n\log\log n})$
a.s. The reverse inequality follows from $\sum_j N_{n,j}=n$, $\sum_j\hat{\rho}_{n,j}=1$
(Lemma~\ref{lem:counts_sum}): $n\hat{\rho}_{n,k}-N_{n,k}=\sum_{j\ne k}(N_{n,j}-n\hat{\rho}_{n,j})$,
which inherits the same one-sided bounds. Combining gives the absolute-value
statements. The last line adds $n(\hat{\rho}_{n,k}-v_k)=O(\sqrt{n\log\log n})$ from
Theorem~\ref{thm:LLN}.
\end{proof}

\section{Joint central limit theorem for the estimators}
\label{sec:norm_clt}

\begin{lemma}[Componentwise joint CLT]\label{lem:componentwise}
\uses{def:estimator, def:counts, def:condA}
Consider any RAR procedure with $N_{n,k}\to\infty$ a.s.\ for all $k$. Under
Condition~\textbf{A}, with $D_n := \mathrm{diag}(\sqrt{N_{n,1}},\dots,\sqrt{N_{n,K}})$,
\[
  D_n(\estParams_n - \Params) \eqdist \mathcal{N}\!\big(0, \mathrm{diag}(V_1,\dots,V_K)\big).
\]
\end{lemma}

\begin{proof}
\uses{lem:Q_martingale, lem:Q_quad_var, thm:mart_clt, cor:mart_clt_componentwise}
The vector martingale $M_{n,k}=\sum_{i\le n} X_{i,k}Y_{i,k}$ (with $Y_{i,k}=\xi_{i,k}-\theta_k$)
has quadratic variation $\langle M_{\cdot,k}\rangle_n = V_k N_{n,k}$ and zero
cross-variation (Lemma~\ref{lem:Q_quad_var}). The Lindeberg condition holds: for
$\varepsilon>0$, $L_{n,k}(\varepsilon)\le \frac1{V_k}\E[Y_{1,k}^2\indic_{\{|Y_{1,k}|>\varepsilon\sqrt{V_k}\sqrt m\}}]$
on $\{N_{n,k}\ge m\}$, which is small for $m$ large by dominated convergence
(using $\E|Y_{1,k}|^2<\infty$), and $\Prob(N_{n,k}<m)\to0$. The martingale CLT
gives $M_n/\sqrt{\langle M\rangle_n}\eqdist \mathcal{N}(0,I_K)$. Finally
$D_n(\estParams_n-\Params) = \mathrm{diag}(\sqrt{V_1},\dots,\sqrt{V_K})\,M_n/\sqrt{\langle M\rangle_n}$
(up to a Bahadur $o_p(1)$ term), giving the stated limit.
\end{proof}

\textbf{Lean remark.} Lemma~\ref{lem:componentwise} uses only $N_{n,k}\to\infty$,
not any property of the limiting proportions. It is the crux of both the
non-sparse and the sparse (Chapter~\ref{chap:forced}) analyses and should be a
standalone theorem depending on a martingale-CLT lemma from Mathlib.

\section{The normality theorem}
\label{sec:norm_main}

\begin{lemma}[CLT for the estimator]\label{lem:clt_theta}
\lean{AlphaRAR.estimator_sqrtN_joint_tendsto_multivariateGaussian}\leanok
\uses{def:estimator, lem:mart_clt_joint_selfnorm, lem:estimator_bahadur}
Suppose $N_{n,k}/n\to v_k$ a.s.\ with $v_k>0$ for every arm $k$. Then, under
Condition~\textbf{A}, $\sqrt n(\estParams_n - \Params)\eqdist\mathcal{N}(0,V)$ with
$V=\mathrm{diag}(V_1/v_1,\dots,V_K/v_K)$.
\end{lemma}

\begin{proof}\leanok
\uses{lem:mart_clt_joint_selfnorm, lem:estimator_bahadur}
Exactly as in Corollary~\ref{cor:mart_clt_componentwise}, but with the deterministic
normalization $\sqrt n$ in place of $D_n=\mathrm{diag}(\sqrt{N_{n,k}})$. By the Bahadur
identity (Lemma~\ref{lem:estimator_bahadur}), coordinate $k$ of
$\sqrt n(\estParams_n-\Params)$ is
$(Q_{n,k}/\sqrt{N_{n,k}})\cdot\sqrt{nN_{n,k}}/(N_{n,k}+1)
 + (\theta_{0,k}-\theta_k)\sqrt n/(N_{n,k}+1)$. The self-normalized martingale vector
$(Q_{n,k}/\sqrt{N_{n,k}})_k$ converges to $\mathcal N(0,\mathrm{diag}(V_k))$
(Lemma~\ref{lem:mart_clt_joint_selfnorm}); the multiplicative factor tends to $1/\sqrt{v_k}$ and
the additive one to $0$ in probability. Multivariate Slutsky then gives the diagonal rescaling
$\mathrm{diag}(1/\sqrt{v_k})$ of the Gaussian, whose covariance is $\mathrm{diag}(V_k/v_k)=V$.
\end{proof}

\begin{lemma}[CLT for the plug-in target]\label{lem:clt_rho}
\lean{AlphaRAR.clt_rho}\leanok
\uses{def:plugin_target}
Under the standing hypotheses, $\sqrt n(\hat{\rho}_n - v)\eqdist\mathcal{N}(0,GVG^\top)$.
\end{lemma}

\begin{proof}\leanok
\uses{lem:clt_theta, lem:rho_rate, lem:taylor_rho}
By the Taylor expansion (as in Lemma~\ref{lem:ell_rho_control}),
$\sqrt n(\hat{\rho}_n-v)=G\sqrt n(\estParams_n-\Params)+O_p(1/\sqrt n)$. Apply
Lemma~\ref{lem:clt_theta} and Slutsky.

\emph{Formalization.} The delta method is assembled in
\texttt{AlphaRAR.clt\_rho}. First, the product-space Slutsky lemma
(\texttt{tendsto\_map\_comp\_of\_tendstoInMeasure\_const}) with $g(x,r)=Gx+r$ combines the estimator
CLT (Lemma~\ref{lem:clt_theta}) with the fact that the first-order remainder
$\sqrt n(\Target(\estParams_n)-\Target(\Params))-G\sqrt n(\estParams_n-\Params)=\sqrt n\,\varphi(\estParams_n-\Params)$
tends to $0$ in probability, and the limiting law is the linear pushforward of the estimator's Gaussian
(\texttt{AlphaRAR.multivariateGaussian\_map\_matrix}, giving the covariance $GVG^\top$). The
remainder-in-probability is discharged by \texttt{MeasureTheory.tendstoInMeasure\_smul\_littleO\_of\_tight}
(an abstract $\varepsilon$--$\delta$ lemma: if $X_n=a_n\cdot S_n$ is tight, $S_n\to0$ in probability and
$\varphi=o(\|\cdot\|)$ at the origin, then $a_n\cdot\varphi(S_n)\to0$ in probability), whose tightness
input is supplied from the estimator CLT via
\texttt{MeasureTheory.tight\_of\_tendsto\_probabilityMeasure} (portmanteau plus continuity from above of
the limit Gaussian over $\{\|x\|\ge j\}\downarrow\emptyset$). The differentiability of $\Target$ at
$\Params$ (Jacobian $G$; Condition~\textbf{B}) and the a.s.\ consistency $\estParams_n\to\Params$
(Theorem~\ref{thm:LLN}) enter as the hypotheses \texttt{hTderiv} and \texttt{hcons}.
\end{proof}

\begin{theorem}[Asymptotic normality]\label{thm:normality}
\uses{def:generic_cond, def:condA, def:condB, def:counts, def:plugin_target, def:estimator}
Under the generic conditions and Conditions~\textbf{A}--\textbf{B}, as
$n\to\infty$:
\begin{enumerate}
  \item[(i)] for every $k$, $|N_{n,k}-n\hat{\rho}_{n,k}|=o_p(\sqrt n)$,
        $|N_{n,k}-n\hat{\rho}_{n,k}|=O(\sqrt{n\log\log n})$ a.s., and
        $N_{n,k}-nv_k=O(\sqrt{n\log\log n})$ a.s.;
  \item[(ii)] $\sqrt n(\estParams_n-\Params)\eqdist\mathcal{N}(0,V)$ and
        \[
          \begin{pmatrix}\sqrt n(N_n/n - v)\\ \sqrt n(\hat{\rho}_n - v)\end{pmatrix}
          \eqdist \mathcal{N}(0,\Omega).
        \]
\end{enumerate}
\end{theorem}

\begin{proof}
\uses{lem:prop_dev, lem:clt_theta, lem:clt_rho}
Part (i) is Lemma~\ref{lem:prop_dev}. For (ii), the marginal CLT for
$\estParams_n$ is Lemma~\ref{lem:clt_theta}. By part (i),
$\sqrt n(N_n/n-\hat{\rho}_n)=o_p(1)$, so $\sqrt n(N_n/n-v)$ and
$\sqrt n(\hat{\rho}_n-v)$ have the same limit $\mathcal{N}(0,GVG^\top)$
(Lemma~\ref{lem:clt_rho}) and are asymptotically equal; the joint vector
$(\sqrt n(N_n/n-v),\sqrt n(\hat{\rho}_n-v))$ therefore equals
$(\sqrt n(\hat{\rho}_n-v),\sqrt n(\hat{\rho}_n-v))+o_p(1)$, which converges to
$\mathcal{N}(0,\Omega)$ by Slutsky.

\emph{Formalization.} The joint statement of part (ii) is
\texttt{AlphaRAR.clt\_joint}: it takes the proportion-deviation fact
$\sqrt n(N_n/n-\hat{\rho}_n)\to_p 0$ of part (i) as a hypothesis and assembles the joint CLT from
\texttt{AlphaRAR.clt\_rho}. The joint vector is the image of $\sqrt n(\hat{\rho}_n-v)$ under the
duplication map $x\mapsto(x,x)$ (the stacked matrix $[I;I]$) up to an $o_p(1)$ remainder, so the
same product-space Slutsky lemma as in Lemma~\ref{lem:clt_rho} applies, and the limiting covariance
$[I;I]\,GVG^\top\,[I;I]^\top$ is the block matrix $\Omega$ via the rectangular Gaussian pushforward
\texttt{AlphaRAR.multivariateGaussian\_map\_matrix}. Part (i)
(Lemma~\ref{lem:prop_dev}, the $o_p(\sqrt n)$ deviation machinery) is not yet formalized, so the
theorem is not marked complete.
\end{proof}

\begin{remark}[Asymptotic efficiency]\label{rem:efficiency}
\uses{thm:normality, prop:exp_fam}
When the arm responses belong to one-parameter exponential families with mean
parameter, Proposition~\ref{prop:exp_fam} gives $V_k = I_k(\theta_k)^{-1}$, so the
asymptotic variance $GVG^\top$ of $N_n/n$ equals the lower bound
$G\,I(\Params)^{-1}G^\top$ of \cite{hu2006asymptotically}, with
$I(\Params)=\mathrm{diag}(v_1 I_1(\theta_1),\dots,v_K I_K(\theta_K))$. Hence every
$\aRTS$ design is asymptotically efficient.
\end{remark}
